\documentclass[11pt]{article}
\usepackage{nopageno}
\setcounter{secnumdepth}{1}
\usepackage[margin=1in]{geometry}
\begin{document}

\pagestyle{empty}

\vspace{-0.5cm}
\subsection{Overview}
\vspace{-0.3cm}
PIs Napolitano (structural vulnerability and community co-design),
Wauthier (SAR processing and uncertainty quantification),
and Busse (hydrological modeling and mass balance),
all at Penn State,
propose a closed-loop cyber-physical system (CPS) for flood duration estimation,
pre-code building vulnerability assessment, and community retrofit decision support
in small riverine communities.
The small, aging communities most repeatedly devastated by riverine flooding lack
the data, software, and technical capacity that existing flood damage frameworks
require.
Community partners in Montpelier, VT have identified inundation
\textit{duration}, not just peak flood stage, as the decision-relevant variable
motivating every modeling constraint in this proposal.

The proposed system has four coupled layers.
The \textbf{physical layer} is the flood-building interaction: pre-code
load-bearing masonry and timber structures absent from every existing vulnerability
framework.
The \textbf{sensing layer} is multi-temporal SAR imagery, open Sentinel-1 C-band
for flood-extent context and commercial ICEYE X-band for building-level duration
and depth at sub-daily cadence.
The \textbf{cyber layer} is a Bayesian fusion framework integrating SAR-derived
observations with hydrological mass balance into probabilistic duration estimates
with propagated uncertainty.
The \textbf{decision layer} is a multi-criteria retrofit allocation framework,
co-developed with community stakeholders, that maps duration outputs to
building-specific damage probabilities, placing community-scale drainage
investments and individual building retrofits on the same probabilistic basis.
A two-mode decision interface delivers the outputs. In Plan/Forecast mode, users
select any flood scenario (historical replay, recurrence interval, or live NWS
forecast) and compare retrofit investments; in Update Model mode, community
members record per-building flood observations that update the Bayesian priors
for future scenarios. Human decisions are the actuators at every timescale,
making the community a functional component of the control loop.

\vspace{-0.5cm}
\subsection{Intellectual Merit}
\vspace{-0.3cm}
The intellectual merit rests on two contributions.
First, we establish that community-accessible sensing, fused under a Bayesian
architecture, can substitute for resource-intensive physics-based simulation in
closed-loop CPS decision systems, a generalizable principle for infrastructure
contexts without hydrodynamic modeling capacity.
Second, we advance a co-design methodology that treats locally held knowledge as
structural input to model calibration and characterizes whether community knowledge
functions as a substitute for or complement to national geospatial data, providing
the first empirical basis for deciding when elicitation improves CPS accuracy.
Two Montpelier flood events (2023, 2024) provide a within-community train-test
validation: the inference chain is calibrated on 2023 and validated against the
held-out 2024 event, directly testing generalizability across event characteristics.
ICEYE commercial SAR provides building-level flood depth and duration at sub-daily
cadence for 271 downtown buildings across both events, establishing the empirical
foundation for fragility calibration and held-out validation.

\vspace{-0.5cm}
\subsection{Broader Impacts}
\vspace{-0.3cm}
The framework enables duration-based damage assessment where no path currently
exists, resolving the allocation problem, community-scale drainage versus
individual building retrofits, that annually re-injures small historic
municipalities.
The decision layer and co-design process are designed for communities without
ongoing research support.
One postdoc and one PhD student receive cross-disciplinary training in SAR
remote sensing, probabilistic structural modeling, hydrological mass balance,
and participatory community research.
The Community-Inspired Research (CIR) co-design methodology, documented as a
transferable protocol, establishes how resident knowledge anchors model calibration
and community decisions feed back into physical infrastructure sensing.
The framework will be released under an MIT license, enabling communities beyond
Montpelier to apply duration-based assessment using freely available national data.

\end{document}
